\section{Rado's graphs and universal functions: proofs}
\label{sec:annex-p03-rado-1964}

\begin{annexitem}{annex:p03.rado64.order-type-boundary}{thm:p03-rado64-order-type-boundary}{The fixed final vertex of an omega-plus-one domain}
\begin{proof}
Let \(h\in F(2,\omega+1,\{0,1\})\). Every domain-preserving
representation injection \(\varphi\) must be strictly increasing:
for \(x<y\), the admissible pair \((x,y)\) requires
\(\varphi(x)\leq\varphi(y)\), and injectivity excludes equality.
Its final value must be \(\varphi(\omega)=\omega\), since a finite
value would leave only finitely many possible images for the infinitely
many predecessors of \(\omega\).

Define \(f\) to be zero on every admissible tuple except that
\(f(\omega,\omega)=1-h(\omega,\omega)\).
The repeated final-maximum pair is in the domain. A representation in
\(h\) would give \(f(\omega,\omega)=h(\omega,\omega)\), contrary
to this definition. Thus \(h\) is not universal. For the countable
initial-ordinal domain, take \(n=0\), \(l=2<\omega\) and two
colours in the first case of \cref{thm:p03-rado64-transfinite}.
There are no GCH hypotheses below index zero.
\end{proof}
\end{annexitem}

\begin{annexitem}{annex:p03.rado64.arity-obstruction}{thm:p03-rado64-arity-obstruction}{Diagonalizing repeated-vertex profiles}
\begin{proof}
Fix any proposed universal \(h\). For each integer \(j\geq0\), set
\(a_j=1-h(j,\ldots,j)\), where this tuple has length \(j+1\).
Define \(f(x_0,\ldots,x_j)=a_j\) on every tuple of that length in the
stated domain. If an injection \(\varphi\) represented \(f\) in \(h\),
put \(j=\varphi(0)\) and apply the representation to the tuple of
\(j+1\) copies of zero. Its image is \(j+1\) copies of \(j\), so it would
give \(a_j=h(j,\ldots,j)=1-a_j\), impossible for a zero-one value.
Repeated coordinates make every tuple used here admissible.
\end{proof}
\end{annexitem}

\begin{annexitem}{annex:p03.rado64.graph-universal}{thm:p03-rado64-graph-universal}{The high-digit embedding}
\begin{proof}
Repeated division with remainder by \(k\geq2\) gives the finite base-\(k\)
expansion of a nonnegative integer. Uniqueness follows by comparing the
least position where two digit strings differ, modulo the next power of \(k\).
These elementary facts justify each digit comparison below.

For a countably infinite input label its vertices by \(\mathbb N_0\).
At stage \(\xi\), the previous images are a finite increasing set.
Choose \(\psi(\xi)\) larger than all of them. The sum below \(2^{\psi(\xi)}\)
uses distinct binary positions, so adding the single high digit introduces
no carry. The resulting \(\varphi(\xi)\) is at least \(2^{\psi(\xi)}\),
strictly larger than all previous images.
For \(x<\xi\), its digit at position \(\varphi(x)\) is exactly the
indicator of \(xE_\Gamma\xi\). Thus both adjacency and nonadjacency
are preserved on each pair. The union of the finite stages is the required
induced embedding.
A finite input can be enlarged by isolated vertices to a countably infinite
one and the embedding restricted back. The empty input has the empty
embedding.
\end{proof}
\end{annexitem}

\begin{annexitem}{annex:p03.rado64.extension}{thm:p03-rado64-extension}{Binary witnesses and complete back-and-forth}
\begin{proof}
Choose integers \(M>\max(U\cup V)\), with the empty-set maximum
condition omitted, and put
\[
 z_M=2^M+\sum_{u\in U}2^u.
\]
Its highest nonzero digit is at \(M\), so \(z_M\) lies above all vertices
in \(U\cup V\). The digits at their positions give exactly the requested
adjacencies. Different sufficiently large \(M\) give different witnesses.

For two countable extension graphs \(G,H\), enumerate both vertex sets.
Starting with any given finite partial induced isomorphism, alternate stages.
At a forward stage choose the first unused vertex \(a\) of \(G\).
Partition the finite current domain into its neighbours and non-neighbours
of \(a\), and transport the partition to \(H\). The extension property
chooses a new vertex outside the current range with exactly this pattern.
Adding this pair preserves an induced partial isomorphism.
At the backward stage apply the same construction to the inverse map and
the first unused vertex of \(H\).
Every vertex eventually appears in the domain and range. The union is a
bijection preserving and reflecting every edge, since any pair occurs
together at some finite stage. This proves uniqueness. With \(G=H\)
and the prescribed initial partial isomorphism, it proves homogeneity.

The existing graph of \cref{thm:p03-fraisse54-graph} has this extension
property, so the constructed isomorphism identifies it with \(R\).
The complete extension-axiom theory and its countable categoricity were
proved there and transfer under the isomorphism. This last step does not
attribute those later model-theoretic formulations to the 1964 article.
\end{proof}
\end{annexitem}

\begin{annexitem}{annex:p03.rado64.amalgams}{thm:p03-rado64-amalgams}{Free graph amalgamation and the forest obstruction}
\begin{proof}
An induced subgraph of a finite simple graph is finite and simple.
Disjoint union supplies joint embedding. For embeddings of a common
finite graph \(A\) into \(B_1,B_2\), take disjoint copies of the parts
outside \(A\), identify exactly the two copies of each element of \(A\),
and keep the edges of both sides. They agree on the identified part because
both maps are induced embeddings. Add no cross-edges.
Each original graph then remains an induced subgraph of a finite simple
graph, proving amalgamation.

In a putative forest amalgam of the two displayed paths, the image of \(x\)
cannot equal the image of \(y\): \(x\) is adjacent to \(b\), while \(y\)
is not. It cannot equal the image of \(z\): \(x\) is adjacent to \(a\),
while \(z\) is not. The embeddings also keep new vertices distinct from
\(a,b\) and keep \(y\ne z\).
Consequently the amalgam contains the five distinct vertices and edges
\(a-x-b-z-y-a\), a cycle. Adding other edges cannot remove this cycle.
A forest amalgam is impossible.
\end{proof}
\end{annexitem}

\begin{annexitem}{annex:p03.rado64.no-local-universal}{thm:p03-rado64-no-local-universal}{The de Bruijn diagonal tree}
\begin{proof}
Label the vertices of the locally finite candidate \(H\) by
\(1,2,\ldots\). For each \(m\geq1\), put \(p_m=k(H,m,m)\).
Finite-step reachability is finite by induction in the number of steps:
the next reachable set is a finite union of finite neighbour sets.

Construct a ray \(x_1-x_2-x_3-\cdots\) and attach \(p_m\) distinct
leaves \(y_{m,1},\ldots,y_{m,p_m}\) to \(x_m\).
All named vertices are distinct. The result \(T\) is countably infinite,
connected, acyclic and locally finite. From \(x_1\), each of these
\(p_m\) leaves and the vertex \(x_{m+1}\) is reachable in exactly \(m\)
steps, so \(k(T,x_1,m)\geq p_m+1\).

A weak embedding preserves walks, and injectivity keeps their endpoints
distinct. Thus it would give
\(k(T,x_1,m)\leq k(H,\varphi(x_1),m)\) for every \(m\).
Set \(m=\varphi(x_1)\), an integer at least one. Then
\[
 p_m<k(T,x_1,m)\leq k(H,m,m)=p_m,
\]
a contradiction. The construction depends on the proposed host, as a
nonuniversality argument must.
\end{proof}
\end{annexitem}

\begin{annexitem}{annex:p03.rado64.combinadic}{thm:p03-rado64-combinadic}{Rado's consecutive-rank calculation}
\begin{proof}
There are finitely many tuples preceding any fixed increasing
\(\ell\)-tuple in colexicographic order, since their largest entry is
bounded by its largest entry. The first tuple is
\((0,1,\ldots,\ell-1)\), of code zero.

Write a tuple's initial consecutive block as
\((a,a+1,\ldots,a+r)\), where \(0\leq r\leq\ell-1\), stopping at the
first gap or at the end. Its immediate successor replaces this block by
\((0,1,\ldots,r-1,a+r+1)\) and leaves later entries unchanged.
The zero prefix contributes zero to the binomial code.
Pascal's identity telescopes to
\[
 \sum_{i=0}^r\binom{a+i}{i+1}
   =\sum_{i=0}^r
       \left(\binom{a+i+1}{i+1}-\binom{a+i}{i}\right)
   =\binom{a+r+1}{r+1}-1.
\]
Thus the code increases by exactly one at each successor.
Colexicographic order is an infinite well-order with finite initial segments:
a nonempty collection has a least largest coordinate, then the same
argument chooses its lower coordinates. It therefore has order type
\(\omega\). Starting at zero and advancing by one proves the asserted
bijection. This is the source's adjacent-tuple argument with the
telescoping identity made explicit.
\end{proof}
\end{annexitem}

\begin{annexitem}{annex:p03.rado64.uniform-complex}{thm:p03-rado64-uniform-complex}{Coding every fixed-size incidence and nonincidence}
\begin{proof}
Label an infinite input by \(\mathbb N_0\).
At stage \(\xi\), for each \(\ell\)-element subset \(A\) of its earlier
vertices, let \(e(A)\) be one or zero according as
\(A\cup\{\xi\}\) is or is not a simplex.
Choose \(\psi(\xi)\) larger than all previous images and all
\(\tau(\varphi[A])\). These are finitely many integers.
Set
\[
 \varphi(\xi)=2^{\psi(\xi)}
   +\sum_{A\in[\xi]^\ell}e(A)\,2^{\tau(\varphi[A])}.
\]
Here the image subset is written in increasing order before applying
\(\tau\). Its distinct codes give distinct digit positions by the lemma,
all below the fresh high digit. Hence the new image exceeds the preceding
ones. Comparing the relevant digit proves equivalence of simplex membership
for every set whose largest input vertex is \(\xi\), including nonmembership.
Every simplex has such a largest vertex, so this completes the embedding.

For \(\ell=0\), choose successive increasing integers of the required
parity, according to membership of each singleton. Both parity classes
are unbounded. Finite inputs are handled by adjoining unused vertices and
restricting the resulting embedding.
\end{proof}
\end{annexitem}

\begin{annexitem}{annex:p03.rado64.mixed-complex}{thm:p03-rado64-mixed-complex}{The all-finite-subsets digit code}
\begin{proof}
On finite subsets the function \(\rho(A)=\sum_{a\in A}2^a\) is injective and onto
\(\mathbb N_0\) by uniqueness and existence of binary expansion,
with the empty set corresponding to zero.

At stage \(\xi\), consider every subset \(A\subseteq\xi\), including
the empty subset. There are exactly \(2^\xi\) such subsets, a finite
number at this countable construction stage.
Let \(e(A)\) record whether \(A\cup\{\xi\}\) is a simplex.
Choose \(\psi(\xi)\) above all previous images and all \(\rho(\varphi[A])\),
and define
\[
 \varphi(\xi)=2^{\psi(\xi)}
       +\sum_{A\subseteq\xi}e(A)\,2^{\rho(\varphi[A])}.
\]
Distinct image subsets have distinct codes; no carry occurs, and every
required incidence or nonincidence is the appropriate digit of the new
vertex. Empty \(A\) gives its singleton membership. Every nonempty finite
set has a largest element and is covered at that stage. This proves
universality, with finite and empty bases handled by restriction.
No condition on the empty simplex has entered.
\end{proof}
\end{annexitem}

\begin{annexitem}{annex:p03.rado64.finite-colours}{thm:p03-rado64-finite-colours}{The base-\(k\) constructions and original digit identities}
\begin{proof}
For fixed \(\ell\geq1\), the source first has an unrestricted digit
function \(f_0(A,t)=d_k(t,\tau(A))\), giving
\[
 t=\sum_{A\in[\mathbb N_0]^\ell}
           k^{\tau(A)}f_0(A,t).
\]
This is the full base-\(k\) representation because the fixed-length \(\tau\)
enumerates every position. The universal function restricts this rule to
\(A<t\), meaning every member of \(A\) is below \(t\).

For the mixed code, the source's representation is
\[
 t=\sum_{A\subseteq\{0,\ldots,t-1\}}
                 k^{\rho(A)} f_k^*(A,t),
\]
where \(f_k^*(A,t)\) denotes the increasing tuple ending at \(t\).
For \(t\geq1\), these positions comprise \(0,\ldots,2^t-1\) and include
every nonzero base-\(k\) digit position of \(t\). The assertion for
\(t=0\) holds with its sole zero-valued empty-prefix term.

For any desired \(f\), construct increasing \(\varphi\) recursively.
At stage \(\xi\), let \(\mathcal A\) be \([\xi]^\ell\) in the fixed-size
case, all subsets of \(\xi\) in the mixed case, and
\(\{\varnothing\}\) in the unary case. Write \(c\) for the corresponding
code \(\tau,\rho\), or the constant zero code.
Choose \(\psi(\xi)>0\) above the previous images and all \(c(\varphi[A])\).
Put
\[
 \varphi(\xi)=k^{\psi(\xi)}
       +\sum_{A\in\mathcal A} f(A,\xi)\,k^{c(\varphi[A])}.
\]
Each coefficient is a digit less than \(k\), and the positions are distinct
and below \(\psi(\xi)\). The lower part is less than \(k^{\psi(\xi)}\), with
no carries; the new image is at least \(k^{\psi(\xi)}>\psi(\xi)\), so it is
strictly above all earlier images.
Digit comparison gives
\(f(A,\xi)=f^*(\varphi[A],\varphi(\xi))\) for every required \(A\).
This proves both parts and the separate unary boundary.
\end{proof}
\end{annexitem}

\begin{annexitem}{annex:p03.rado64.restriction}{thm:p03-rado64-restriction}{Colour retraction and arity restriction}
\begin{proof}
Choose \(s_0\in S_0\) and a retraction \(r:S\to S_0\), taking every
element outside \(S_0\) to \(s_0\).
For a universal \(f^*\) on the larger domain, define the proposed smaller
one by restricting arities and then applying \(r\).
Given \(f_0\in F(l_0,m,S_0)\), extend it to \(F(l,m,S)\) by the constant
\(s_0\) on the other arities. A representing injection for this extension
is also domain-preserving on the smaller arities. At each such tuple
the larger universal value equals \(f_0\), which is fixed by \(r\).
This proves the reduction, including the empty arity domain.
\end{proof}
\end{annexitem}

\begin{annexitem}{annex:p03.rado64.transfinite}{thm:p03-rado64-transfinite}{Template enumeration in the regular and singular cases}
\begin{proof}
Identify cardinals with initial ordinals in this proof and write
\(\kappa=\aleph_n\). If \(l=0\), the unique empty function is universal.
Assume \(l>0\), and fix a default colour in the nonempty set \(S\).

\paragraph{Cardinal arithmetic used.}
In ZFC, \(\theta\cdot\theta=\theta\) for every infinite cardinal \(\theta\).
One elementary proof is induction on infinite cardinals: well-order pairs
by their maximum coordinate and then lexicographically within each layer.
Each proper initial segment at a pair of maximum \(\alpha<\theta\)
has cardinal below \(\theta\), using the result for the smaller infinite
cardinal \(\max(\aleph_0,|\alpha|)\); finite initial layers are immediate,
including the base \(\theta=\aleph_0\).
If the order type exceeded \(\theta\), its element at position \(\theta\)
would have an initial segment of cardinal \(\theta\), a contradiction.
The reverse inequality follows from a coordinate injection.
This also gives products and unions of at most \(\theta\) sets of size
at most \(\theta\) size at most \(\theta\).
Exponential identities used below follow from identifying functions on a
product with families of functions, for example
\((2^\alpha)^\beta=2^{\alpha\cdot\beta}\).

An infinite successor cardinal \(\theta^+\) is regular: a cofinal family
of at most \(\theta\) smaller ordinals would cover it by at most \(\theta\)
sets each of size at most \(\theta\), contrary to its cardinality.
A limit cardinal with GCH below it is strong limit, since every smaller
infinite cardinal is some \(\aleph_\nu\) and
\(2^{\aleph_\nu}=\aleph_{\nu+1}\) remains below the limit.
The countable cardinal is strong limit for finite arguments as well.

For a regular strong-limit \(\kappa\), one has
\(\kappa^\mu=\kappa\) for \(0<\mu<\kappa\).
Every map \(\mu\to\kappa\) has bounded image by regularity.
For each bound, the number of maps is below \(\kappa\):
for infinite parameters, bound it by \(2^\theta<\kappa\) for some
infinite \(\theta<\kappa\) containing both parameters.
Finite parameters give the same conclusion directly.
Taking \(\kappa\) possible bounds gives at most \(\kappa\) maps;
constant maps supply the reverse inequality.

We also need the following consequence of GCH below an infinite \(\theta\):
\[
 \theta^\mu\leq\theta\qquad(\mu<\operatorname{cf}\theta).
\]
For a limit \(\theta\), each such map has bounded image and the preceding
strong-limit bound, summed over \(\theta\) bounds, proves it.
For successor \(\theta=\zeta^+\), GCH gives \(\theta=2^\zeta\);
then for \(0<\mu<\theta\), \(\zeta\mu=\zeta\) and
\(\theta^\mu=2^{\zeta\mu}=\theta\).
For \(\theta=\aleph_0\), \(\mu\) is finite and finite products suffice.
The exponent zero is immediate.

\paragraph{Counting templates in cases (i) and (ii).}
For a bound \(\eta<\kappa\), a template has domain
\(\bigsqcup_{\lambda<l}D_\eta^\lambda\).
In case (i), \(\kappa\) is regular strong limit, with \(l<\kappa\).
This domain has cardinal \(\delta<\kappa\).
For the countable case both the bound and the arity set are finite;
for uncountable \(\kappa\), use an infinite
\(\theta<\kappa\) bounding \(|\eta|,|l|\), so the domain size is at most
\(\theta\,2^\theta<\kappa\).
There are at most \(\kappa^\delta\leq\kappa\) templates at that bound.
With \(\kappa\) bounds, the complete set of templates has size at most
\(\kappa\).

In case (ii), put \(\theta=\aleph_r\), so
\(\kappa=\theta^+=2^\theta\) and \(l<\operatorname{cf}\theta\).
A bound \(\eta<\kappa\) has cardinal at most \(\theta\);
the just-proved power bound shows that the template domain has size at most
\(\theta\). The number of templates there is at most
\(\kappa^\theta=(2^\theta)^\theta=2^\theta=\kappa\).
Again there are at most \(\kappa\) templates altogether.

\paragraph{The regular construction, including repeated maxima.}
Enumerate all templates with repetition as
\((\eta_\nu,h_\nu)_{\nu<\kappa}\).
Partition \(\kappa\) into \(\kappa\) disjoint sets \(P_\nu\), each of
cardinality \(\kappa\), using a bijection with \(\kappa\times\kappa\).
Each \(P_\nu\) is unbounded, since a bounded subset of the initial ordinal
\(\kappa\) has cardinal less than \(\kappa\).
Assign every \(z\in P_\nu\) the template \((\eta_\nu,h_\nu)\), and define
\(f^*\) by \cref{def:p03-rado64-marked-template}.
This is one fixed function on the whole required domain.

Given \(f\), suppose increasing \(\varphi\) is defined below \(\xi<\kappa\).
Regularity of \(\kappa\) supplies a bound \(\eta<\kappa\) strictly above
all previous images. Define a template \(h\) at \(\eta\).
For every \(\lambda<l\) and prefix \((x_i)_{i<\lambda}\) with \(x_i\leq\xi\),
replace \(x_i<\xi\) by \(\varphi(x_i)\), and \(x_i=\xi\) by \(*\);
assign this encoded prefix the value \(f((x_i)_{i<\lambda},\xi)\).
The encoding is injective on prefixes of each length, so there is no
conflict, even when an earlier coordinate equals \(\xi\) many times.
Give all other template arguments the default colour.
Choose \(\nu\) representing this template and
\(\varphi(\xi)\in P_\nu\) with \(\varphi(\xi)\geq\eta\).
This extends the increasing map.

For any tuple with final coordinate \(\xi\), its earlier strict-predecessor
images lie below \(\eta\), while every occurrence of \(\varphi(\xi)\)
is sent to the marker. The definition of \(f^*\) therefore gives exactly
the prescribed value of \(f\). Transfinite recursion proves universality
in both regular cases. No value of the new image was needed when the
template was chosen.

\paragraph{The singular block construction.}
In case (iii), \(\kappa\) is singular. Its index \(n\) is a nonzero limit
ordinal, since successor cardinals and \(\aleph_0\) are regular.
Choose \(m_0<n\) with
\(|S|,|l|\leq\aleph_{m_0}\).
For every \(m_0\leq m<n\), let
\(\kappa_m=\aleph_{m+2}\), and form a block \(X_m\) of order type
\(\omega_{m+2}\). Order
\[
 X=\sum_{m_0\leq m<n}X_m
\]
lexicographically by block and then within the block.
Each \(\kappa_m<\kappa\) is regular.
The initial segment before a point of \(X_m\) has cardinal at most
\(\aleph_{m+1}\): earlier blocks number at most \(|m|\) and each has
cardinal at most \(\aleph_{m+1}\), and a proper initial segment within
\(X_m\) has at most that cardinal.
All proper point-initial segments therefore have cardinal below \(\kappa\),
whereas the block cardinalities are cofinal in \(\kappa\).
The order type of \(X\) is exactly \(\omega_n\); if it exceeded this initial
ordinal, the point at position \(\omega_n\) would contradict the
initial-segment bound.

For \(\eta\in X_m\), its predecessor set with the marker has cardinal
at most \(\aleph_{m+1}\).
Since \(|l|,|S|\leq\aleph_m\), the number of arguments of a template
is at most
\[
 \aleph_m(\aleph_{m+1})^{\aleph_m}
 =\aleph_{m+1},
\]
using \(\aleph_{m+1}=2^{\aleph_m}\).
The number of templates at a fixed bound is at most
\[
 (\aleph_m)^{\aleph_{m+1}}
 \leq2^{\aleph_{m+1}}=\aleph_{m+2}=\kappa_m.
\]
The same bound holds after including all \(\kappa_m\) choices of \(\eta\)
inside \(X_m\). Enumerate them with repetition by \(\mu<\kappa_m\).
Partition \(X_m\) into \(\kappa_m\) disjoint subsets \(P_{m,\mu}\),
each of cardinal \(\kappa_m\) and unbounded in its block.
For a last argument \(z\in P_{m,\mu}\), define \(f^*\) using that
block's indexed marked template.

Construct increasing \(\varphi:X\to X\) preserving each block.
At a stage \(\xi\in X_m\), earlier images in earlier blocks are below
the entire current block. There are fewer than \(\kappa_m\) previous
images within \(X_m\), so regularity of \(\kappa_m\) bounds them by
some \(\eta\in X_m\).
Use precisely the same marked-prefix prescription as in the regular case,
now on the whole predecessor set of \(\eta\) in \(X\).
Choose its indexed cell \(P_{m,\mu}\) and choose the new image there at
least \(\eta\).
The recursion remains in its assigned block and strictly increases.

For any argument sequence in the source domain, use its final maximum
\(\xi\). Strict predecessors have previously assigned images, and all
coordinates equal to \(\xi\) are handled by the marker.
The same verification gives equality of \(f\) and \(f^*\) on its image.
Transport through the order isomorphism \(X\cong\omega_n\) proves the
third case. This retains the source's strict colour bound and all
arity and cardinal-arithmetic hypotheses.
\end{proof}
\end{annexitem}
